% =============================================================================
% Project : THE_PHILOSOPHY_OF_VIABILITY
% Book    : Book I — Ontology of Viability
% Book ID : VS.B01
%
% File    : CHAPTER_01_IDENTITY/sections/
%           section_01_ontological_problem_of_identity.tex
%
% Section:
%   The Ontological Problem of Identity
%
% Purpose:
%   Establishes the ontological problem of identity, clarifies
%   its relation to being and determination, distinguishes
%   identity from related ontological concepts, and defines
%   the scope and limits of the present investigation.
%
% =============================================================================


\section{The Ontological Problem of Identity}
\label{sec:ontological-problem-of-identity}


% -----------------------------------------------------------------------------
% 1. THE QUESTION OF WHAT AN ENTITY IS
% -----------------------------------------------------------------------------

\subsection{The Question of What an Entity Is}

The question of identity occupies a foundational
position in ontology because it concerns what it
means for an existent to be the existent that it is.

The question is not merely how an entity differs
from other entities, how it can be characterized,
or which properties can be attributed to it.
It concerns the ontological constitution of an
existent and the status of its identity within
the actuality of its being.

To ask what an entity is therefore involves more
than establishing a list of its characteristics.
It requires examining the relation between the
entity's existence and the identity through which
it is the entity that it is.

This problem cannot be resolved by treating identity
as an additional property acquired by an entity
after its existence has already been established.
Such an approach would presuppose an entity whose
existence could be understood independently of
its being anything determinate.

Nor can identity be treated as an independent
ontological condition that precedes the existence
of the entity to which it belongs. Both approaches
separate what this investigation regards as
ontologically inseparable.

The starting point of the present inquiry is
therefore the ontological status of identity itself.
Identity must be examined as a dimension of the
actuality of what exists, rather than as a merely
descriptive characterization imposed upon it.

This formulation establishes the central problem
of the section: what is the ontological significance
of being the entity that an existent is, and how
does this identity relate to the determinations
through which its being is constituted and articulated?


% -----------------------------------------------------------------------------
% 2. IDENTITY AS AN ONTOLOGICAL PROBLEM
% -----------------------------------------------------------------------------

\subsection{Identity as an Ontological Problem Rather Than a Descriptive Problem}

Identity is frequently approached through properties,
characteristics, identifying marks, or configurations
of features. Such approaches can be useful for
characterizing entities and examining their
distinguishing features. Nevertheless, they do
not necessarily establish the ontological status
of the identity of the entities concerned.

A description may specify characteristics of
an entity without establishing whether those
characteristics are constitutive of its identity.
A distinguishing feature may differentiate one
entity from another without explaining the
ontological basis of that distinction. A
configuration of properties may characterize
an existent without exhausting what that
existent is.

The distinction between identity and its
characterization is therefore indispensable.

The identity of an entity must not be equated
with the set of statements, properties, or
descriptive features associated with it.
Such features may be relevant to its ontological
constitution, but their descriptive availability
does not establish their constitutive status.

Likewise, a characteristic may distinguish
an entity within a particular context without
being sufficient to establish its identity
in the full ontological sense.

The problem becomes more significant when
different kinds of determination are considered
together. An entity may involve structural,
relational, and material determinations. It may
participate in relations, undergo transformations,
and exhibit different manifestations. Yet the
presence of these dimensions does not, by itself,
establish how they contribute to the constitution
of its identity.

The ontological investigation must consequently
distinguish three interconnected questions:

\begin{enumerate}
    \item What is the ontological status of identity
    as a dimension of an existent?

    \item Through which determinations and relations
    is the identity of an existent constituted
    and articulated?

    \item Under what conditions can an existent
    undergo alterations in its determinations
    without necessarily losing or replacing
    its identity?
\end{enumerate}

These questions are related, but they are not
interchangeable. The first concerns the ontological
problem of identity itself. The second concerns
the relation between identity and determination.
The third introduces the problem of identity
in relation to continuity and transformation.

A satisfactory ontology of identity must preserve
these distinctions rather than treating the
description or characterization of an entity
as a sufficient explanation of its being.


% -----------------------------------------------------------------------------
% 3. THE PROBLEM OF IDENTITY AND BEING
% -----------------------------------------------------------------------------

\subsection{The Problem of Identity and Being}
\label{subsec:identity-and-being-problem}

The ontological status of identity immediately
raises the question of its relation to being.

If identity is treated as a property that an
already existing entity subsequently acquires,
the account implicitly assumes that an entity
can first exist independently of being anything
determinate.

Such an assumption creates a fundamental
ontological difficulty. An entirely identity-free
existent would lack any established basis for
being the existent that it is. The attempt to
describe what exists prior to its identity
would therefore risk presupposing the very
determination whose absence it seeks to establish.

The position developed in this work is that
being and identity are ontologically inseparable.
To exist is not to exist first as an identity-free
substrate and subsequently become an entity
through the addition of identity. Rather,
the actuality of an existent and its being
the existent that it is belong to the same
ontological actuality.

This inseparability does not eliminate the
analytical distinction between being and identity.

Being concerns the actuality of what exists,
whereas identity concerns the existent as
the existent that it is. These formulations
indicate distinguishable dimensions of one
ontological actuality; they do not establish
two independently existing components.

Three claims must nevertheless be kept distinct.

First, being and identity are analytically
distinguishable: each can be examined under
a different conceptual aspect.

Second, being and identity are ontologically
inseparable: neither is treated here as an
independently existing component that can
be detached from the other.

Third, neither is assigned ontological priority
over the other. Their analytical distinction
does not establish an order of ontological
dependence, and their inseparability does
not, by itself, demonstrate the claim of
non-priority. Non-priority is maintained
as a distinct foundational commitment.

Accordingly, the relation between being and
identity is neither one of temporal succession
nor one of external addition. Identity is not
imposed upon being, and being is not an
identity-neutral foundation that must receive
its identity from elsewhere.

This position constitutes a foundational
ontological commitment of the Viability
Philosophy System. Its implications and
conceptual limits will be examined more
precisely in Section~\ref{sec:being-and-identity}.


% -----------------------------------------------------------------------------
% 4. THE PROBLEM OF IDENTITY AND DETERMINATION
% -----------------------------------------------------------------------------

\subsection{The Problem of Identity and Determination}

The inseparability of being and identity
leads to a further question: what is the
relation between identity and determination?

An existent is not an entirely indeterminate
actuality. Its being involves a determinate
ontological character. Nevertheless, it would
be premature to conclude that every determination
associated with an existent contributes to
its identity in the same way or possesses
the same ontological status.

Some determinations may be constitutive of
the existent, whereas others may arise through
its relations, conditions, manifestations,
or participation in particular ontological
domains. Some may persist through transformations,
while others may vary or cease to obtain.

The ontological significance of these differences
cannot be assumed in advance.

The problem is therefore not simply whether
an existent is determinate, but how its
determinations relate to its identity.

This distinction prevents two opposite reductions.

On the one hand, identity cannot be separated
from determination as though an existent could
possess an actuality entirely independent of
any determinate character.

On the other hand, identity cannot be reduced
to an exhaustive inventory of all properties,
relations, or conditions associated with
an existent. The presence of a determination
does not automatically establish that it is
constitutive of identity, and the accumulation
of determinations does not, by itself,
explain their ontological unity or significance.

An adequate account must therefore investigate
which determinations are constitutive, how
they are related, and whether their alteration
is compatible with the continued identity
of the existent.

These questions will be developed in
Section~\ref{sec:identity-and-determination}.
At the present stage, it is sufficient to
establish that determination is indispensable
to the investigation of identity without
assuming that all determinations possess
the same ontological status.


% -----------------------------------------------------------------------------
% 5. THE PROBLEM OF IDENTITY AND DIFFERENCE
% -----------------------------------------------------------------------------

\subsection{The Problem of Identity and Difference}

Identity also raises the problem of difference.

To consider an existent as the existent that
it is involves distinguishing its identity
from the identity of what it is not. However,
the existence of a difference between entities
does not, by itself, explain what constitutes
the identity of either entity.

Difference and identity are therefore related
but non-equivalent ontological concerns.

A distinction may involve a particular feature,
relation, structural configuration, or material
determination. Yet the fact that a characteristic
differentiates two entities does not establish
that this characteristic constitutes the
identity of either one.

Distinguishability cannot automatically be
equated with ontological constitution.

This becomes particularly important when
considering entities that share several
determinations while remaining distinct,
as well as entities whose determinations
may vary under different conditions.

The problem of difference must therefore
be addressed without reducing identity
to the accumulation of distinguishing
properties.

The investigation will return to this issue
through the analysis of unity, distinction,
individuation, numerical identity, and
signature. At this stage, the essential
commitment is that difference may articulate
distinctions between entities without
necessarily providing a complete account
of what each entity is.


% -----------------------------------------------------------------------------
% 6. THE PROBLEM OF IDENTITY THROUGH CHANGE
% -----------------------------------------------------------------------------

\subsection{The Problem of Identity Through Change}

A further difficulty emerges when identity
is considered in relation to change.

Entities may undergo transformations in
their conditions, relations, structures,
and manifestations. Their determinations
may vary, their configurations may be
reorganized, and their relations with
other entities may be altered.

These possibilities raise a fundamental
ontological question: under what conditions
does an existent remain the existent that
it is through transformation?

The question cannot be answered simply by
asserting that identity requires the persistence
of all properties. Such a position would make
identity incompatible with any alteration
in the determinations of an existent.

Conversely, it cannot be assumed that every
existent preserves its identity regardless
of the transformations it undergoes.

Both positions require arguments that have
not yet been established.

The problem of identity through change therefore
requires an account of the relationship between
constitutive determinations and changing
conditions. It also requires distinguishing
the persistence of an existent from the
persistence of particular characteristics
associated with it.

The relation between identity and continuity
must consequently be investigated without
presupposing that continuity is equivalent
to the uninterrupted persistence of every
determination.

Likewise, the possibility of transformation
does not by itself establish either the
preservation or the replacement of identity.

The detailed examination of continuity,
persistence, and transformation requires
further conceptual and ontological analysis.
The present section establishes these issues
as problems to be investigated rather than
assuming a definitive criterion in advance.


% -----------------------------------------------------------------------------
% 7. THE DISTINCTION BETWEEN IDENTITY AND RELATED CONCEPTS
% -----------------------------------------------------------------------------

\subsection{The Distinction Between Identity and Its Related Concepts}

The ontological problem of identity cannot
be adequately formulated without distinguishing
identity from several closely related concepts.

Unity concerns the question of how an existent
constitutes a unity. It must not be assumed
that unity is equivalent to singularity,
that every unity is an aggregate, or that
unity can be explained solely through composition.

Individuation concerns the question of how
an existent is determined or distinguished
as this existent rather than another. Its
relation to identity requires independent
analysis; individuation must not simply be
substituted for identity.

Numerical identity concerns whether an existent
at different times, positions, or conditions
is one and the same existent. This question
is indispensable to the analysis of persistence
and continuity, but numerical identity alone
does not provide a complete account of the
ontological constitution of an existent.

Invariance concerns the preservation of
specified features or relations under
specified transformations and with respect
to a specified domain. It may be relevant
to the examination of identity, but the
presence of an invariant does not by itself
establish identity, and invariance is not
equivalent to persistence or continuity.

Signature concerns a structured characterization
associated with an existent. A signature may
help articulate or distinguish an existent,
but its descriptive or identifying usefulness
does not establish that it constitutes the
identity of that existent.

Determination concerns the ways in which
an existent is determinate. It is indispensable
to the inquiry into identity, but identity
cannot be equated with an exhaustive inventory
of determinations.

Expression and manifestation introduce further
questions concerning the articulation and
manifestation of ontological determinations.
Their relation to identity must be examined
without treating either concept as synonymous
with identity or presuming that either exhausts
the constitution of an existent.

These concepts are related to identity because
they address different aspects of the constitution,
differentiation, characterization, persistence,
or manifestation of existents. Nevertheless,
their relevance does not establish their
equivalence.

The principal distinctions required at this
stage are summarized in Table~\ref{tab:identity-conceptual-distinctions}.
The table provides a conceptual orientation
for the chapter; it does not claim to offer
final or exhaustive definitions of every
listed concept.

% =============================================================================
% Project : THE_PHILOSOPHY_OF_VIABILITY
% Book    : Book I — Ontology of Viability
% Book ID : VS.B01
%
% Asset   : Table
% ID      : VS.B01.P01.C01.S01.SS07.P001.TL01
% Name    : identity_conceptual_distinctions
%
% File    :
%   4_ASSETS/tables/
%   VS.B01.P01.C01.S01.SS07.P001.TL01.identity_conceptual_distinctions.tex
%
% Referenced from:
%   Chapter 01 — Section 01 — Subsection 07
%
% =============================================================================

\begin{table}[htbp]
\centering
\caption{Conceptual Distinctions Between Identity and Related Ontological Concepts}
\label{tab:identity-conceptual-distinctions}

\renewcommand{\arraystretch}{1.35}
\small

\begin{tabular}{|p{0.23\textwidth}|p{0.67\textwidth}|}
\hline
\textbf{Concept} & \textbf{Distinction from Identity} \\
\hline

Being
&
Concerns the actuality of what exists. It is
analytically distinguishable from identity but
is not treated as ontologically separable from it.
Neither is assigned ontological priority over
the other.
\\
\hline

Determination
&
Concerns the determinate character of an existent.
Determinations may be relevant to the constitution
of identity, but identity is not defined as the
sum or inventory of all determinations.
\\
\hline

Unity
&
Concerns the constitution or coherence of an
existent as a unity. Unity is not automatically
equivalent to singularity, aggregate, composition,
or identity.
\\
\hline

Individuation
&
Concerns the determination or differentiation
through which an existent is distinguished as
this existent. Its relation to identity requires
analysis and cannot be assumed to be equivalent
to identity.
\\
\hline

Numerical Identity
&
Concerns whether an existent considered across
different times, positions, or conditions is
one and the same existent. This question does
not, by itself, exhaust the ontological constitution
of that existent.
\\
\hline

Difference
&
Concerns distinction or non-identity between
entities or determinations. A difference may
distinguish entities without necessarily
establishing what constitutes the identity
of either one.
\\
\hline

Invariance
&
Concerns the preservation of specified features
or relations under specified transformations
and with respect to a specified domain. An
invariant does not, by itself, establish
identity, persistence, or continuity.
\\
\hline

Signature
&
Concerns a structured characterization associated
with an existent. A signature may contribute
to its characterization or differentiation
without necessarily constituting its identity.
\\
\hline

Expression
&
Concerns the articulation of determinations
or aspects of an existent. It must not be
treated as synonymous with identity or as
an exhaustive account of its constitution.
\\
\hline

Manifestation
&
Concerns how an existent or its determinations
are manifested. A manifestation is not,
without further argument, identical to the
identity of that which is manifested.
\\
\hline

Continuity
&
Concerns a relation of connectedness or
continuation across conditions or stages.
It must not be equated in advance with
identity preservation or with the persistence
of every determination.
\\
\hline

\end{tabular}

\end{table}

% =============================================================================
% END OF TABLE
% =============================================================================

The distinctions in Table~\ref{tab:identity-conceptual-distinctions}
are not intended to establish a complete theory
of every listed concept.

Their purpose is to prevent conceptual substitution:
a concept related to identity must not be treated
as identical to identity without an independent
ontological argument.

Nor should conceptual distinction be interpreted
as proof of ontological independence. The fact
that two concepts can be analytically distinguished
does not establish that the corresponding
ontological dimensions exist independently
of one another.


% -----------------------------------------------------------------------------
% 8. THE SCOPE OF THE ONTOLOGICAL INVESTIGATION
% -----------------------------------------------------------------------------

\subsection{The Scope of the Ontological Investigation}

The scope of the present investigation
must be clearly delimited.

First, identity is examined as an ontological
dimension of what exists. Its existence
and constitution are not made dependent upon
the presence or activity of an external
identifying subject.

Second, the investigation does not assume
that the identity of an existent can be
exhaustively captured through a set of
properties or distinguishing characteristics.
Such characteristics may be relevant to
its constitution, but their constitutive
significance must be established rather
than presumed.

Third, the investigation does not assume
that an existent necessarily possesses
only one identity. The possibility of
multiple identities within one existent,
including the possibility that different
respects or ontological levels may be relevant
to different aspects of identity, remains
an open question. Neither a plurality of
manifestations nor a plurality of components
is, by itself, evidence of a plurality
of identities.

Fourth, the investigation does not equate
identity preservation with the persistence
of all properties or the absence of change.
The relationship between identity and
transformation requires an account of
constitutive determinations and the
conditions under which their alteration
may be compatible with identity.

Fifth, the investigation does not assume
that identity is reducible to composition,
structure, relation, materiality, or
manifestation. These dimensions may be
relevant to the ontological constitution
of an existent, but their precise roles
must be established through the subsequent
development of the ontology.

Finally, the present section does not
attempt to provide a complete account
of the organization of the realms of being,
the ontological status of light, or the
dynamics through which determinations
are articulated and transformed.

These questions require the development
of the subsequent investigations within
the architecture of the Ontology of Viability.

These limitations are not arbitrary restrictions.
They preserve the distinction between the
foundational problem of identity and the
additional ontological questions required
for its more comprehensive explanation.


% -----------------------------------------------------------------------------
% 9. THE CENTRAL ONTOLOGICAL PROBLEM
% -----------------------------------------------------------------------------

\subsection{The Central Ontological Problem}

The preceding considerations allow the
central problem of this section to be
formulated more precisely.

The question is not merely how an existent
can be distinguished from other existents,
or how its characteristics can be described.
The question is what constitutes an existent
as the existent that it is, how this identity
belongs to its ontological actuality, and
how it relates to the determinations and
transformations through which its being
is articulated.

This problem contains several interconnected
dimensions:

\begin{enumerate}
    \item The analytical distinction and
    ontological inseparability of being
    and identity, together with the distinct
    commitment that neither has ontological
    priority over the other.

    \item The relation between identity
    and the determinations that may be
    constitutive of an existent.

    \item The distinction between identity,
    unity, distinction, and individuation.

    \item The relation between identity
    and difference.

    \item The distinction between identity
    and its descriptive or characterizing
    dimensions, particularly signature.

    \item The relationship between identity,
    numerical identity, continuity, and
    transformation.

    \item The possibility of plurality
    of identities within a single existent,
    without presupposing either unity or
    plurality as the necessary outcome.
\end{enumerate}

These dimensions cannot be resolved through
a single definition or by reducing identity
to one of its associated concepts.

They require a systematic ontological
investigation in which the relations among
them are established without prematurely
collapsing their distinctions.

The purpose of this chapter is therefore
to establish the conceptual and ontological
foundations necessary for such an investigation.
It begins with the inseparability of being
and identity, proceeds to the relationship
between identity and determination, examines
unity, distinction, and individuation, and
subsequently investigates numerical identity,
signature, and the principle of non-reduction.

The chapter does not presume that all these
relations have already been resolved. Rather,
it establishes the problem-space within which
their ontological significance can be
rigorously examined.


% -----------------------------------------------------------------------------
% 10. CONCLUDING FORMULATION
% -----------------------------------------------------------------------------

\subsection{Concluding Formulation}

Identity is not an additional property of
an existent, nor is it merely a descriptive
representation of its characteristics or
a criterion through which it can be
distinguished from other existents.

It is an ontological concern inseparable
from the actuality of being.

To investigate identity is therefore to
investigate what it means for an existent
to be the existent that it is, how this
identity belongs to its ontological
constitution, and how it relates to the
plurality of determinations, relations,
and transformations through which its
being is articulated.

The central position established here
distinguishes three claims: being and
identity are analytically distinguishable;
they are ontologically inseparable; and
neither is assigned ontological priority
over the other. These claims are related,
but none is treated as an automatic
demonstration of either of the others.

This position provides the foundation
for the subsequent examination of identity.
It also establishes a necessary restriction:
no related concept may replace identity
without an independent ontological argument,
and no analytical distinction may be taken
as sufficient proof of ontological independence.

The ontological problem of identity is thus
established as the starting point for the
systematic investigation undertaken in
the following sections.


% =============================================================================
% END OF SECTION
% =============================================================================